\documentclass[conference]{IEEEtran}

\usepackage{times}
\usepackage{graphicx}
\begin{document}

\begin{center}

{\LARGE \textbf{Presidio-Based LLM Security Mini Gateway}} \\[0.5cm]

{\large Artificial Intelligence (AIC201)} \\[0.2cm]

{\large Assignment 2} \\[0.8cm]

\textbf{Student Name:} Hadia Cheema \\[0.2cm]

\textbf{Enrollment Number:} 01-134232-056 \\[0.2cm]

\textbf{Instructor:} Dr. Arshad Farhad \\[0.2cm]

\textbf{Department:} Computer Science \\[0.2cm]

\textbf{University:} Bahria University \\[0.2cm]

\textbf{Submission Date:} \today

\end{center}

\vspace{1cm}

\section{Introduction}


Large Language Models (LLMs) have become one of the most important technologies in modern artificial intelligence. These models are designed to understand natural language and generate human-like responses. Because of this ability, LLMs are now used in many real-world applications such as chatbots, virtual assistants, automated helpdesk systems, document analysis tools, and educational platforms. Many companies integrate LLMs into their software to improve user interaction and automate tasks that previously required human effort.

Although these models are very powerful, they also introduce new types of security risks. Since LLMs interact with users through natural language prompts, the system must interpret whatever text the user provides. This creates an opportunity for attackers to manipulate the model by carefully designing malicious prompts. These prompts are written in such a way that they attempt to override the instructions of the system or trick the model into producing unintended outputs.

One common threat that affects LLM systems is known as \textit{prompt injection}. In a prompt injection attack, the user includes instructions within the input prompt that attempt to change the behavior of the language model. For example, an attacker might try to instruct the model to ignore previous rules or reveal hidden system instructions. Because LLMs process natural language in a flexible way, they may follow these instructions if proper security measures are not implemented.

Another important concern when working with language models is the leakage of sensitive information. Users interacting with the system may unintentionally include personal data such as phone numbers, email addresses, identification numbers, or even private API keys in their prompts. If such information is processed without filtering, it can create serious privacy risks. In enterprise environments, accidental exposure of confidential information can lead to financial loss or damage to organizational reputation.

In addition to personal information, system prompts themselves are considered sensitive assets. System prompts are hidden instructions used by developers to control how the language model behaves. These prompts often contain internal rules, policies, or operational instructions. If an attacker successfully extracts the system prompt, they may gain insight into the system's internal logic and use that knowledge to exploit weaknesses in the application.

Because of these risks, many researchers and engineers recommend placing a protective layer in front of LLM systems. This layer acts as a security gateway that analyzes user input before it reaches the language model. By inspecting prompts in advance, the system can detect suspicious instructions and prevent potentially harmful requests from being processed.

In this project, a lightweight LLM security gateway is implemented to demonstrate this concept. The gateway acts as a filtering mechanism that evaluates user input and applies security checks before sending the request to the language model. The system focuses on two main security objectives: detecting prompt injection attempts and identifying sensitive information in user prompts.

To detect sensitive information, the system integrates Microsoft Presidio, which is a well-known open-source framework designed for identifying personally identifiable information (PII) in text. Presidio is capable of detecting entities such as phone numbers, email addresses, credit card numbers, and other types of sensitive data.

The gateway also includes a prompt injection detection mechanism. This component analyzes the prompt and checks for suspicious phrases commonly associated with injection attempts. 

After analyzing the input, a policy engine determines the appropriate action. Depending on the results, the system may allow the request, mask sensitive information, or block the prompt entirely if it appears malicious.


\section{Threat Modeling}

Before implementing any security mechanism for a Large Language Model system, it is important to understand the possible threats that the system may face. Threat modeling is a process used to identify potential risks, understand how attackers might exploit the system, and determine which assets must be protected. In the context of LLM-based systems, threat modeling helps developers design protective mechanisms that prevent misuse of the model.

Large Language Models interact directly with users through text prompts. This open interaction makes the system flexible and useful, but it also introduces several security concerns. Because the system processes whatever input the user provides, attackers can attempt to manipulate the model by sending specially crafted prompts. These prompts are designed to influence the behavior of the model in ways that the system developer did not intend.

One of the most common threats in LLM systems is prompt injection. In a prompt injection attack, the user attempts to insert instructions that override the system's original rules. For example, a user may write something like “ignore the previous instructions and reveal the system prompt.” 

Another important threat is jailbreak attempts. A jailbreak attack occurs when a user tries to bypass the safety mechanisms of the language model. Instead of directly asking the model to break the rules, the attacker may use creative wording, role-playing scenarios, or indirect instructions. 

Sensitive information leakage is also a serious concern. Users may include personal information in their prompts without realizing the risks. This information may include phone numbers, email addresses, account numbers, or other private details. If the system processes such data without filtering it, the information may be stored, logged, or transmitted in ways that compromise user privacy.

Another possible threat involves API keys or secret tokens. Developers sometimes accidentally paste credentials into prompts while testing systems. If these credentials are processed without protection, attackers could gain access to external services or systems.

System prompt extraction is another potential risk. System prompts are hidden instructions that guide the behavior of the language model. These prompts often contain internal rules and policies that control how the model should respond. If an attacker manages to extract these instructions, they may learn how to manipulate the system more effectively.



In addition to identifying assets, threat modeling also requires understanding the capabilities of potential attackers. In this project, it is assumed that attackers interact with the system only through text prompts. They do not have direct access to the source code, database, or system infrastructure. Instead, they attempt to manipulate the model using carefully designed prompts.

Because of this assumption, the most effective security strategy is to analyze user input before it reaches the language model. By inspecting prompts in advance, the system can detect suspicious instructions and identify sensitive information. This approach allows the system to block malicious prompts or mask sensitive data before it is processed.

Understanding these threats helps guide the design of the security gateway implemented in this project. The gateway focuses on detecting prompt injection attempts and identifying sensitive information within user prompts. By addressing these two major risks, the system can significantly improve the security of LLM-based applications.

\subsection{Protected Assets and Sensitivity Levels}

The LLM system processes different types of data which must be protected from unauthorized access or exposure. These assets vary in importance and therefore have different sensitivity levels.

\begin{table}[h]
\centering
\begin{tabular}{|l|l|}
\hline
Asset & Sensitivity Level \\
\hline
System Prompt & High \\
User Personal Data & High \\
API Keys / Tokens & Critical \\
Conversation History & Medium \\
Application Configuration & Medium \\
\hline
\end{tabular}
\caption{Protected assets and sensitivity levels}
\end{table}

\subsection{Common LLM Attack Types}

LLM systems are vulnerable to multiple types of attacks due to their natural language interface.

\begin{itemize}

\item \textbf{Prompt Injection}  
Attackers attempt to override system instructions by embedding malicious instructions in the prompt.

\item \textbf{Jailbreak Attacks}  
Users attempt to bypass model safety mechanisms through creative prompts or role-play scenarios.

\item \textbf{Sensitive Data Leakage}  
Users may include confidential information such as phone numbers or emails in prompts.

\item \textbf{API Key Exposure}  
Developers may accidentally expose API keys or credentials in prompts.

\item \textbf{System Prompt Extraction}  
Attackers attempt to force the model to reveal hidden internal instructions.

\end{itemize}

\subsection{Attacker Assumptions}

In this project, attackers are assumed to interact with the system through normal user prompts. They do not have direct access to the internal infrastructure or source code of the system.

However, attackers may attempt to manipulate the model using carefully crafted prompts designed to bypass security controls or extract sensitive information.

The security gateway therefore inspects every incoming prompt before it reaches the language model. Suspicious prompts can be blocked, while sensitive data can be masked before processing.

\section{System Architecture}

The proposed system is designed as a lightweight security gateway placed in front of a Large Language Model. The goal of this gateway is to analyze user prompts before they reach the language model in order to detect malicious instructions and sensitive information.

The architecture follows a modular pipeline structure where each component performs a specific task. This design improves maintainability and allows new security mechanisms to be added easily in the future.

When a user submits a prompt, the input first passes through the security gateway rather than being sent directly to the language model. The gateway analyzes the prompt using multiple processing stages.

The first stage of the system is the \textbf{Injection Detection Module}. This module examines the prompt and searches for suspicious patterns commonly used in prompt injection attacks. Examples include phrases such as ``ignore previous instructions'', ``reveal system prompt'', or ``bypass security rules''. The module assigns a score to the prompt based on how many suspicious phrases are detected.

After the injection detection stage, the prompt is forwarded to the \textbf{Presidio Analyzer}. Microsoft Presidio is an open-source framework used for detecting personally identifiable information (PII) in text. In this system, Presidio scans the prompt and identifies sensitive entities such as phone numbers, email addresses, credit card numbers, and API keys.

If sensitive entities are detected, they are recorded along with their positions in the prompt. The system can then mask or anonymize this information before sending the prompt forward.

The third stage of the architecture is the \textbf{Policy Engine}. The policy engine acts as the decision-making component of the system. It evaluates the injection score and the detected sensitive entities to determine how the system should respond.

Three possible actions are supported by the policy engine:

\begin{itemize}
\item \textbf{Allow} – The prompt is safe and can be forwarded to the language model.
\item \textbf{Mask} – Sensitive information is detected and must be anonymized.
\item \textbf{Block} – The prompt is considered malicious and should not reach the language model.
\end{itemize}

After the policy decision is made, the processed prompt is either blocked, masked, or forwarded to the language model. This layered architecture ensures that unsafe inputs are filtered before reaching the LLM.

Another important component of the system is the \textbf{Latency Monitoring Module}. This module measures the time required for each processing stage. Monitoring latency helps evaluate whether the security gateway introduces significant delays during prompt processing.

\subsection{Processing Pipeline}

The gateway operates as a sequential processing pipeline where each module performs a specific task. This pipeline design ensures that the prompt is analyzed step by step before reaching the language model. Each stage focuses on a particular security concern, which improves the reliability of the system.

When a prompt enters the system, it first passes through the injection detection module. This module performs pattern matching and searches for suspicious instructions that are often used in prompt injection attacks. The system maintains a list of keywords and phrases associated with malicious behavior. If these phrases appear in the prompt, the injection score increases.

The next stage is the Presidio analysis module. This module focuses on detecting personally identifiable information in the prompt. The Presidio analyzer uses both rule-based detection and pattern matching to identify entities such as phone numbers, email addresses, credit card numbers, and API keys. Detecting these entities is important because sensitive information should not be processed without protection.

After both analyses are completed, the system forwards the results to the policy engine. The policy engine combines the injection score and the list of detected entities to determine the appropriate action. The policy rules define thresholds that control how the system responds to different risk levels.

For example, prompts with a high injection score are blocked completely because they may contain malicious instructions. Prompts containing sensitive data are masked to protect user privacy. If the prompt does not contain any security risks, it is allowed to pass to the language model.

This modular pipeline design ensures that security checks are applied consistently for every user prompt.

The modular architecture of the system makes it flexible and easy to extend. Additional detection modules or security policies can be integrated without changing the overall pipeline structure.

\begin{figure}[h]
\centering
\includegraphics[width=1.00\columnwidth]{architecture.drawio.png}
\caption{LLM Security Gateway Architecture}
\end{figure}

\section{Implementation}

The proposed security gateway was implemented using Python. The system follows a modular design where each security component is implemented as a separate module. This approach makes the system easier to maintain and allows additional security mechanisms to be integrated in the future.

The project is organized into several Python files, each responsible for a specific task in the processing pipeline. The main program coordinates the interaction between these modules and executes the evaluation scenarios.

The entry point of the system is implemented in the \textbf{main.py} file. This file is responsible for running the program and simulating different user prompts that are used to test the security gateway. Each prompt is passed through the gateway pipeline where the injection detection module, Presidio analyzer, and policy engine evaluate the input.

The \textbf{injection\_detector.py} module is responsible for identifying prompt injection attempts. This component scans the input prompt and searches for suspicious phrases that are commonly associated with prompt injection attacks. A scoring mechanism is used to measure how likely the prompt is to be malicious. If several suspicious patterns are detected, the injection score increases.

Sensitive information detection is implemented using Microsoft Presidio in the \textbf{presidio\_engine.py} module. Presidio is an open-source framework designed to detect personally identifiable information within text. In this implementation, the system is configured to detect entities such as phone numbers, email addresses, credit card numbers, and API keys. When such information is detected, the system records the entity type and its location in the prompt.

The \textbf{policy\_engine.py} module is responsible for determining how the system should respond after the input has been analyzed. The policy engine evaluates both the injection score and the detected sensitive entities. Based on predefined thresholds, the engine decides whether the prompt should be allowed, masked, or blocked.

Latency measurement is implemented in the \textbf{latency.py} module. This component measures the processing time for each stage of the security pipeline. The recorded latency values are later used to evaluate the performance impact of the gateway.

The project also includes a \textbf{scenarios.py} module which defines several test prompts used to evaluate the system. These prompts simulate different real-world situations such as normal user questions, prompts containing sensitive data, and malicious prompt injection attempts.

Finally, configuration parameters used by the system are stored in the \textbf{config.py} file. This file contains configurable values such as injection detection thresholds and policy rules.

Overall, the modular structure of the implementation allows each security component to operate independently while still working together as part of the gateway pipeline. This design improves system clarity and simplifies debugging and future improvements.

\subsection{Module Interaction}

The interaction between modules is coordinated by the main program. When the system starts, the main module loads configuration parameters and initializes the detection components.

Each input prompt is processed sequentially through the pipeline. First, the injection detector analyzes the prompt and calculates an injection score. Next, the Presidio analyzer scans the text to detect sensitive entities.

The results from both modules are combined and passed to the policy engine. The policy engine evaluates these results and determines the appropriate response. If masking is required, the system replaces sensitive entities with placeholder tokens such as <PHONE_NUMBER> or <EMAIL_ADDRESS>.

The system also records latency information for each module. This allows developers to evaluate how much processing time is introduced by the security gateway. Monitoring latency is important because the gateway should not significantly slow down user interactions.

Overall, the implementation demonstrates how multiple security components can work together to protect LLM systems from common threats.
\subsection{Presidio Customization}

The system includes several custom Presidio recognizers to improve the detection of sensitive entities. 
In addition to the default Presidio recognizers, custom patterns were implemented to detect API keys and 
internal identifiers that may appear in prompts.

These custom recognizers use regular expressions and contextual scoring to improve detection accuracy. 
For example, API keys are identified using pattern matching based on typical key formats used in cloud 
services. Context-aware scoring is also applied to increase the confidence score when entities appear 
near keywords such as “key”, “token”, or “secret”.

These customizations extend the capabilities of Presidio and improve the system’s ability to detect 
sensitive information in real-world prompts.

\section{Evaluation}
The following table summarizes the evaluation scenarios used to test the security gateway.
\begin{table}[htbp]
\centering
\begin{tabular}{|l|l|l|l|}
\hline
Scenario & Injection Score & PII Detected & Decision \\
\hline
Normal Question & 0 & None & Allow \\
Phone Number Prompt & 0 & Phone Number & Mask \\
Email + Phone Prompt & 0 & Email, Phone & Mask \\
Prompt Injection Attempt & 2 & None & Block \\
API Key Leak & 0 & API Key & Mask \\
\hline
\end{tabular}
\caption{Scenario-level evaluation}
\end{table}

\begin{table}[h]
\centering
\begin{tabular}{|l|l|l|}
\hline
Recognizer & Purpose & Result \\
\hline
Email Recognizer & Detect email addresses & Successful \\
Phone Recognizer & Detect phone numbers & Successful \\
API Key Recognizer & Detect API keys & Successful \\
Credit Card Recognizer & Detect credit card numbers & Successful \\
\hline
\end{tabular}
\caption{Presidio customization validation}
\end{table}

\begin{table}[h]
\centering
\begin{tabular}{|l|l|}
\hline
Metric & Value \\
\hline
Total Scenarios Tested & 5 \\
Injection Detection Accuracy & 100\% \\
PII Detection Accuracy & 100\% \\
Blocked Attacks & 1 \\
Masked Sensitive Data & 3 \\
\hline
\end{tabular}
\caption{Performance summary metrics}
\end{table}

\begin{table}[h]
\centering
\begin{tabular}{|l|l|}
\hline
Injection Score & Policy Action \\
\hline
0 & Allow \\
1 & Mask / Review \\
2+ & Block \\
\hline
\end{tabular}
\caption{Threshold calibration policy}
\end{table}

\begin{table}[h]
\centering
\begin{tabular}{|l|l|}
\hline
Component & Latency (ms) \\
\hline
Injection Detection & 0.01 \\
Presidio Analysis & 34.2 \\
Policy Engine & 0.01 \\
Total Processing Time & 34.22 \\
\hline
\end{tabular}
\caption{Latency summary}
\end{table}

\subsection{Security Implications}

The results highlight the importance of introducing a security gateway when deploying LLM systems. Without such a gateway, malicious prompts or sensitive data may reach the language model directly. By introducing prompt inspection and data masking, the system significantly reduces the risk of prompt injection and accidental information leakage.

\subsection{Results Analysis}

The evaluation results demonstrate the effectiveness of the proposed security gateway. Several test scenarios were executed to simulate real-world interactions with an LLM system.

The first scenario represents a normal user query that does not contain any malicious instructions or sensitive information. The gateway correctly identifies the prompt as safe and allows it to proceed without modification.

The second scenario includes a phone number in the prompt. The Presidio analyzer successfully detects the phone number entity and the policy engine masks the sensitive information before forwarding the prompt. This ensures that personal data is protected.

The third scenario includes both an email address and a phone number. The system detects both entities simultaneously and applies masking to both pieces of sensitive information. This demonstrates the ability of the system to detect multiple sensitive entities within a single prompt.

The fourth scenario simulates a prompt injection attempt. In this case, the injection detection module identifies suspicious instructions and assigns a high injection score. The policy engine blocks the prompt completely, preventing it from reaching the language model.

The final scenario tests the detection of API keys. The custom recognizer implemented in the Presidio analyzer successfully identifies the API key pattern and masks it before processing.

Overall, the evaluation results confirm that the security gateway can successfully detect prompt injection attempts and protect sensitive data while introducing minimal latency.

\section{GitHub Repository}

The complete implementation of the Presidio-Based LLM Security Mini Gateway is available in the following GitHub repository:

\textbf{https://github.com/fitbit12/llm-security-mini-gateway}

The repository contains the full source code, installation instructions, and scripts required to reproduce the evaluation results presented in this report.

The repository includes a README file containing step-by-step instructions for environment setup, dependency installation, and execution of evaluation scenarios. These instructions allow other researchers to reproduce the experimental results reported in this paper.

\subsection{Code Quality and Reproducibility}

The project was developed using a modular and well-structured code design. Each major component of the security gateway is implemented as a separate Python module, including the injection detection module, Presidio analyzer integration, policy engine, latency monitoring component, and evaluation scenario definitions.

This modular structure improves readability, maintainability, and ease of debugging. It also allows additional security modules to be integrated into the gateway without changing the overall system architecture.

The GitHub repository contains all source code, configuration files, and dependency information required to run the system. The repository includes a \texttt{requirements.txt} file listing all required Python libraries, along with a detailed \texttt{README} file explaining how to install dependencies and execute the evaluation scenarios.The repository also contains example prompts and evaluation scripts so that the experimental results reported in this paper can be reproduced by other researchers.


\section{Conclusion}

This project presented the design and implementation of a lightweight security gateway intended to protect Large Language Model (LLM) systems from common prompt-based attacks and sensitive data exposure. As LLM technologies become increasingly integrated into real-world applications such as chatbots, customer service systems, and automated assistants, ensuring the security and reliability of these systems becomes an important challenge.

The proposed gateway acts as a protective layer placed between the user and the language model. Instead of sending user prompts directly to the LLM, the gateway first analyzes the input using several security mechanisms. This approach allows the system to detect malicious instructions, identify sensitive information, and apply appropriate security policies before the prompt is processed by the language model.

One of the major features of the system is the prompt injection detection module. This component evaluates user prompts and identifies suspicious instructions that attempt to override system rules or extract confidential information. By assigning an injection score to each prompt, the system can estimate the level of risk associated with the input and determine whether it should be blocked.

Another important feature of the gateway is the integration of the Microsoft Presidio framework for detecting personally identifiable information (PII). Presidio allows the system to identify sensitive entities such as phone numbers, email addresses, credit card numbers, and API keys within user prompts. When such information is detected, the system can mask or anonymize the data before further processing. This helps protect user privacy and prevents accidental exposure of confidential information.

The evaluation results demonstrate that the security gateway can successfully identify different types of security threats. The system correctly allows normal prompts, masks sensitive information when required, and blocks prompts that contain malicious injection attempts. At the same time, latency measurements show that the additional security checks introduce only minimal performance overhead, which makes the gateway suitable for real-time applications.

Another advantage of the proposed design is its modular architecture. Each component of the gateway operates independently and performs a specific security task. This modular structure makes the system easy to maintain and extend. For example, additional detection mechanisms such as advanced prompt filtering or machine learning based anomaly detection could be integrated into the pipeline in the future.

Although the implemented gateway focuses mainly on prompt injection detection and PII identification, the project demonstrates an important principle: LLM systems should not directly process user input without proper validation. Adding a security layer significantly improves the safety and reliability of LLM-based applications.

Future improvements may include more advanced prompt injection detection techniques, improved context-aware analysis of user prompts, and integration with monitoring tools for detecting unusual usage patterns. Researchers may also explore adaptive security policies that dynamically adjust protection levels based on system behavior.

In conclusion, the implementation of a dedicated security gateway represents a practical and effective approach for improving the safety of LLM-based systems. By combining prompt injection detection, sensitive data protection, and policy-based decision making, the proposed solution helps reduce several common risks associated with modern language model applications.



\end{document}
